\section{A uniform local solution at an integer resonance}
\label{sec:resonance}

At a regular singular point the two usual power-series solutions can collapse
into one. The Gauss hypergeometric report gives a finite recurrence in this
case and leaves a joint coefficient bound and convergence modulus for variable
resonance order to be supplied. We construct both, and determine exactly when
a logarithm is forced. The same construction applies to other scalar
second-order equations once one regular solution and a Wronskian density are
given. The source discussions occur in the resonance paragraphs of sections 5
and 10 of the repository's \texttt{hypergeometric\_ipra\_2026-11-16} report.

\subsection{Local integration with a variable pole order}
Use Gaussian rationals with $|x+iy|_1=|x|+|y|$. This norm is
submultiplicative. For a conditional object theorem the stated bounds on the input coefficient programs are supplied as iPRA derivations or retained as hypotheses. The coefficient of $z^n$ in a product is a sum over
$i+j=n$. A Laurent series here has an explicitly bounded finite principal
part. The arithmetic conclusions below are identities and error inequalities
for finite coefficient and value observations. The language of solution
germs describes the analytic interpretation of those conclusions.

Let $m$ be any natural number. Suppose two fixed PR procedures
(with numerical parameters) present $u=\sum u_nz^n$ and $g=\sum g_nz^n$,
with $u_0=g_0=1$ and a supplied integer $D\ge1$. Assume the arithmetical
hypotheses $|u_n|_1,|g_n|_1\le D^n$ for every $n$. In the Gauss
application below these inequalities follow from the recurrence. Define
\[
 w_0=1,\qquad w_n=-\sum_{i=1}^n u_iw_{n-i}\quad(n>0),\qquad
 h_n=\sum_{i+j+k=n}g_iw_jw_k,\qquad K=6D .
\]
Then $uw=1$ and $u^2h=g$ coefficient by coefficient. Inductively
$|w_n|_1\le(2D)^n$ follows by dividing the recurrence bound by $(2D)^n$ and using $\sum_{i=1}^n 2^{-i}\le1$. Each triple contributes at most
$(2D)^n$, and there are at most $3^n$ triples. Each triple is encoded
injectively by a word of $i$ zeros, $j$ ones and $k$ twos, including
the empty word. Thus $|h_n|_1\le K^n$.

Fix a nonzero Gaussian rational centre $z_*$ with $K|z_*|_1<1/4$.
Use rational samples $z$ with $K|z|_1\le1/4$ and
$|t|_1\le1/2$, where $t=(z-z_*)/z_*$. Set
\[
 L(z)=\sum_{j\ge1}(-1)^{j+1}t^j/j .
\]
Here $L(z_*)=0$ and $L'(z)=1/z$ because the derivative sums to
$1/[z_*(1+t)]$ by the finite geometric identity. The chart guard
guarantees $z\ne0$ since $t=-1$ fails it. Define
\begin{align}
 F(z)&=\sum_{n<m}\frac{h_n}{n-m}z^{n-m}+h_m L(z)+T(z), &
 T(z)&=\sum_{j\ge1}\frac{h_{m+j}}{j}z^j, \nonumber\\
 v(z)&=u(z)F(z), & F'(z)&=z^{-m-1}h(z). \label{eq:frob-pri}
\end{align}
The three terms of $F$ integrate the powers below $-1$, the power $-1$,
and the nonnegative powers. The finite sum is empty when $m=0$.

For $0\le d\le2$, omitting terms $j\ge J\ge 2$ in $T^{(d)}$
has error at most
\begin{equation}
                   4^{d+1}K^{m+d}2^{-J}.                   \label{eq:frob-tail}
\end{equation}
Indeed the coefficient magnitudes of $T$ are bounded by $K^{m+j}$.
The $j$th differentiated term on the chart is bounded by
$4^d K^{m+d}j^d4^{-j}$. For $j\ge2$, $j^d\le2^{j+1}$, by checking
$j=2,3$ and using $(j+1)^2\le2j^2$ for $j\ge3$. Sum the geometric
majorant to obtain \eqref{eq:frob-tail}. The same estimate with $m=0$
applies to $u,w,g,h$. For error $2^{-p}$ take
$J=2+p+4^{d+1}K^{m+d}$ since $C\le2^C$ for integral $C\ge1$.
The log tail is at most $2^{1-J}$ and its first two derivatives are the
rational functions $1/z,-1/z^2$.

Derivative assertions can themselves be read as finite remainder inequalities.
For a polynomial partial sum of a series with $|a_j|_1\le AK^j$, subtract
its value and linear term at $z$. By the binomial theorem the remainder at
$z+s$ has norm at most
\[
 |s|_1^2 AK^2\sum_{j\ge2}\binom j2 2^{-(j-2)}
 =8AK^2 |s|_1^2\qquad (K(|z|_1+|s|_1)\le1/2).
\]
For finite sums the bound $8$ follows by multiplying three finite geometric
sums of ratio $1/2$. Their triple coefficients are $\binom j2$ at
degree $j-2$ up to the truncation. For the derivative polynomial the
remainder is at most $48AK^3|s|_1^2$ by
$j\binom{j-1}{2}=3\binom j3$ and four geometric sums. Applying
\eqref{eq:frob-tail} at the two observed points proves the derivative
identities to arbitrary precision. One may enlarge a boundary chart
slightly within the open disk of convergence, replacing $1/4$ by any
rational ratio strictly below one and using its geometric tail. For $L$
apply the same polynomial proof with $|t|_1+|s/z_*|_1\le3/4$.
In particular no theorem on differentiating a general infinite limit
is needed.

Individual cutoffs give a cutoff for $v$ and each derivative. Evaluate
negative powers at rational $z$ exactly. A magnitude bound for each series
is one plus the norm of its precision-zero approximation. Propagate bounds
for sums and products by the triangle and product inequalities. To request
precision $p$ at a sum request $p+1$ at both children. At a product with
operand bounds $A,B$ request $p+\lceil A+B+2\rceil$ at both children,
since the product error is at most $(A+B+1)\epsilon$ for input
errors at most $\epsilon\le1$. Apply this rule also to the product
formulas for $v'$ and $v''$. It accounts for small denominators from the
principal part.

\paragraph{Local reduction of order.}
Suppose $u$ satisfies $y''+P y'+Qy=0$ on the chart, and
$W=z^{-m-1}g$ satisfies $W'=-PW$. Then $v$ in
\eqref{eq:frob-pri} satisfies the same equation and
$uv'-u'v=u^2F'=W$. Indeed $(u^2F')'=-Pu^2F'$ gives
$u^2F''+(2uu'+Pu^2)F'=0$. Expand the equation for $uF$, use the
equation for $u$, and cancel the unit $u$. The inverse exists since $uw=1$.
Also $|u-1|_1,|g-1|_1\le1/23$ by $D|z|_1\le1/24$.
In particular the inverse Wronskian is $z^{m+1}/g$, where $1/g$ is
computed by the same inverse-prefix rule. At an ordinary centre where
$P,Q$ have convergent Taylor data, the two solutions span the Taylor-solution
germs. Solve the two initial-value equations using this inverse Wronskian.
The difference has zero constant and linear coefficients. At order $j$ the
equation determines its next coefficient by a convolution of earlier ones
divided by $(j+2)(j+1)$, so every coefficient vanishes. Recentring a convergent
series on a smaller disk uses the binomial expansion. At parent degree $j$
the sum of norms of all translated terms is at most
$|a_j|_1(|z_*|_1+|s|_1)^j$, giving a geometric bound for the reordered
tails. Rational negative powers are handled by the local geometric inverse.

A Laurent derivative has zero coefficient at $z^{-1}$ because only the
derivative of a constant could contribute. If a Laurent solution with
a finite pole had normalized Wronskian $W$ with $u$, its quotient by $u$
would have derivative $z^{-m-1}h$, whose residue is $h_m$. Conversely,
if $h_m=0$ the displayed formula gives such a Laurent solution. This
obstruction is an identity of formal Laurent coefficients at zero.
The chart gives the local analytic realization.

\subsection{The Gauss equation and its exact logarithmic coefficient}
Take $a,b\in\mathbb Q(i)$, $c=m+1$, and put
\[
 \mathcal L y=z(1-z)y''+[c-(a+b+1)z]y'-aby.
\]
Starting at $u_0=g_0=1$, compute
\[
 u_{n+1}=\frac{(n+a)(n+b)}{(n+1)(n+m+1)}u_n,\qquad
 g_{n+1}=\frac{n-q}{n+1}g_n,\qquad q=m-a-b .
\]
Comparison of coefficients gives $\mathcal L u=0$ and
$(1-z)g'=-qg$. Write $P=[c-(a+b+1)z]/[z(1-z)]$.
The logarithmic derivative of $z^{-m-1}g$ is
$-(m+1)/z-q/(1-z)=-P$. The integer
\[
 D=\left\lceil\max\{1,\max(1,|a|_1)\max(1,|b|_1),|q|_1\}\right\rceil
\]
works for both series since $n+|a|_1\le(n+1)\max(1,|a|_1)$,
and the same inequality holds for $b$ and $q$. Thus the two local
solutions have a single PR precision schedule uniform in $m,a,b$.

Write $(x)_m=\prod_{j<m}(x+j)$. The logarithm in $v$ has multiplier $h_m u$,
where
\begin{equation}
             h_m=\frac{(-1)^m(a-m)_m(b-m)_m}{(m!)^2}.        \label{eq:frob-res}
\end{equation}
For $m=0$ this gives $h_0=1$. For $m>0$ separate the logarithmic term and
write the Laurent part of $v$ as $z^{-m}\sum_{j\ge0}d_jz^j$. Its leading
coefficient is $d_0=-1/m$. The logarithmic source has no powers below
$z^{-1}$ since
\[
 \mathcal L(u L)=2(1-z)u'+[m/z-(a+b)]u .
\]
For $0\le j<m-1$ the coefficient of $z^{j-m}$ therefore gives
\[
 (j+1-m)(j+1)d_{j+1}=(j+a-m)(j+b-m)d_j.
\]
At $z^{-1}$ the multiplier of $d_m$ is zero and the remaining equality is
$m h_m=(a-1)(b-1)d_{m-1}$. Multiplying through the previous $m-1$ steps
gives \eqref{eq:frob-res} because the denominator product is
$(-1)^{m-1}((m-1)!)^2$, with the other two factors of $m$ from $d_0$ and
the last equation. This includes $m=1$ via the empty product.
Hence for $m>0$ a second Laurent solution exists exactly when
$a$ or $b$ belongs to $\{1,\dots,m\}$, a decidable condition over Gaussian
rationals. In the reverse implication for an arbitrary independent
meromorphic solution, its Wronskian divided by $W$ has derivative zero.
Coefficientwise only the constant survives, and independence makes the
constant nonzero (apartness in the analytic reading). Rescaling would
therefore give a normalized solution.

For $c=1-m$ the substitution $y=z^mY$ gives the Gauss equation for $Y$
with parameters $a+m,b+m,1+m$. The extra coefficient of $z^{-1}Y$
is $m(m-1+c)=0$, the derivative coefficient becomes
$c+2m-(a+b+2m+1)z$, and the constant is $-(a+m)(b+m)$.
Consequently the construction covers every integer $c$ at zero.
Replacing $z$ by $1-z$ changes $c$ to $a+b+1-c$ and gives the analogous
local basis at one whenever $c-a-b$ is integral. The constants matching a
solution from a distant base point and the global Schwarz mapping problem
remain separate questions.

\paragraph{A resonant example.}
For $a=b=1/2$, $m=2$ one has $u_1=1/12$, $g=1-z$, $h_0=1$,
$h_1=-7/6$, and $h_2=9/64$. Hence the normalized second solution is
\[
 v=u\left(-\frac1{2z^2}+\frac7{6z}
             +\frac9{64}L(z)+\sum_{j\ge1}h_{2+j}z^j/j\right).
\]
Here $D=1$, $K=6$ and $z_*=1/48$ satisfies the chart guard. For
$a=1,b=5/3,m=3$, the coefficient $h_3$ vanishes and the second solution
is meromorphic at the singular point.

\paragraph{Object derivation and Lean completion.}
Code the Gaussian-rational pairs with positive denominators and clear them at
comparisons. Compute $(u_i,g_i,w_i,h_i)_{i\le N}$ by appending to bounded
tables. A QF induction assertion is the Boolean fold checking the defining
identities and the bounds for every index in the current prefix. Apply
the same finite convolution calculation to these prefixes to obtain $uw=1$,
$u^2h=g$ and the differential identities, retaining the needed coefficients
(e.g.\ through $m+j+2$ for Laurent powers through $z^j$ and two derivatives).
For \eqref{eq:frob-res} maintain the cross-multiplied recurrence product
through a prefix of length at most $m-1$. For \eqref{eq:frob-tail} take
arbitrary finite endpoints $J\le M$, bound the finite sum through $M$ by
the displayed geometric quantity, then universally close in $M,J,p$.
Use the same finite-endpoint convention for products by bounding the partial
convolution tail at a strict radius by grouping triples by total degree,
using the bound $(6D)^j$ above. These inductions use integer equalities,
cross-multiplied inequalities or bounded Boolean checks. The universal input
bounds on the two coefficient families stay in the context of induction.

The accompanying \texttt{lean/Resonance.lean} implements the real-rational
coefficient tables and product formula, proves the size of the inverse table,
the order-zero coefficient, the bounded zero test for the product formula,
and a precision inequality, and kernel-checks the
two examples. For the object theorem first lift these definitions to pairs of
rational codes including the signed coefficients of the finite principal part.
Next derive the prefix estimates and identities just listed. Lastly use the
error circuits to construct the chart's first two derivative names. The
repository's \texttt{HypergeometricOdeFormal} contains the formal Wronskian
derivative identity, and \texttt{ComplexLogarithmName} constructs a local
logarithm name. The latter leaves the derivative law open. The finite
geometric identity and difference estimate above supply its mathematical
proof after identifying the two local power-series observations. A direct
implementation can use the displayed chart throughout. The supplied Lean
file covers finite lemmas and examples. The full object-language derivation
remains to be built. The independent exact tests in
\texttt{code/resonance.py} compare the convolution and product answers and
check the coefficients of the inhomogeneous Laurent recurrence.
